\documentclass[10pt,a4paper]{article}
\usepackage{amsmath,amssymb,amsthm,graphicx}
\usepackage[margin=25mm]{geometry}
\usepackage{mathtools,tikz-cd}
\usepackage[hidelinks]{hyperref}
\newtheorem{theo}{Theorem}[section]
\newtheorem{prop}[theo]{Proposition}
\newtheorem{corollary}[theo]{Corollary}
\theoremstyle{definition}
\newtheorem{defi}[theo]{Definition}
\numberwithin{equation}{section}
\numberwithin{figure}{section}
\newcommand{\bn}{\mathbf{n}}
\newcommand{\bu}{\mathbf{u}}
\newcommand{\bm}{\mathbf{m}}
\newcommand{\mcM}{\mathcal{M}}
\newcommand{\mcL}{\mathcal{L}}
\newcommand{\mfl}{\mathfrak{l}}
\newcommand{\mfL}{\mathfrak{L}}
\newcommand{\bbN}{\mathbb{N}}
\newcommand{\bbZ}{\mathbb{Z}}
\newcommand{\genn}{\mathrm{gen}}
\DeclareMathOperator{\CDi}{\mathsf{CD}}
\DeclareMathOperator{\Nov}{\mathsf{Nov}}
\DeclareMathOperator{\Mag}{\mathsf{Mag}}
\DeclareMathOperator{\Vect}{Vect}
\renewcommand*{\to}{\mathchoice{\longrightarrow}{\rightarrow}{\rightarrow}{\rightarrow}}
\let\oldmapsto\mapsto
\renewcommand*{\mapsto}{\mathchoice{\longmapsto}{\oldmapsto}{\oldmapsto}{\oldmapsto}}
\newcommand*{\twoheadto}{\mathbin{\mathrlap{\to}\mkern -2mu\to}}
\title{Free multi-Novikov algebras as canonical multidifferential subalgebras (Theorem 3.3)}
\author{Yvain Bruned, Vladimir Dotsenko}
\date{doi:10.5802/ahl.275; CC BY 4.0}
\begin{document}
\maketitle
\noindent\textit{Editorial note.} Native statement and proof are unchanged. Selected bibliography reconstructed only from the same official published PDF.\par\medskip
\noindent\textit{Editorial scope.} Author-material preparation only; one Theorem 3.3, no mathematical repair or formal admission.\par\medskip
\noindent\textit{Allowed standard prerequisites (editorial).}\begin{itemize}
\item Knuth's rotation correspondence. [Knu97], The Art of Computer Programming, Vol. 1, Fundamental Algorithms, 3rd ed. (1997), Chapter 2.3. Standard planar binary-tree/rooted-tree correspondence; actual map, labeling rule and both native diagrams retained at lines 430–492.
\item Free commutative differential algebra and graded dimensions. Definition 3.1, Proposition 3.2, and the finite-set/finer-grading footnote of this source; [Bur06] Section 2.1 provides comparative context. Polynomial realization is stated in full, commuting derivations checked in Proposition 2.1; finite-dimensional surjection/dimension comparison is elementary linear algebra.
\item Earlier single-operation Novikov realization. [DL02] Theorem 7.8; [BN85], [Osb92]; [BCCH21] Proposition 4.21. Earlier realization is acknowledged strategy/context, not a substitute for current proof. All new multi-operation normal-form, spanning and bijection arguments are retained.
\item General multi-index setting. [BL24], published bibliography key BL23. Background motivation for Section 2 variables and operations; target Theorem 3.3 has its own complete source-local algebraic proof.
\end{itemize}\medskip
\noindent\textit{Source and counted unit.} Yvain Bruned and Vladimir Dotsenko, \emph{Novikov algebras and multi-indices in regularity structures}, Annales Henri Lebesgue 9 (2026), 791--814, \url{https://doi.org/10.5802/ahl.275}. Published under CC BY 4.0: \url{https://creativecommons.org/licenses/by/4.0/}.
\par\noindent The counted unit is the whole Theorem 3.3 (free multi-Novikov realization). The canonical-structure proof, populated-monomial and dimension bounds, ordered and normal reductions, and both native rotation-correspondence diagrams are support for that one unit. Sections 2 and 3 below preserve the original numbering; no helper proposition is separately counted.
\par\noindent\textit{Editorial boundary.} The original author mathematics is retained, including its notation and wording. The journal class and branding are replaced by a standard wrapper. The seven cited bibliography entries are transcribed from the published bibliography, pp.812--813, with original printed labels; no missing BibTeX file is guessed. Knuth's rotation correspondence and elementary free differential-algebra/graded dimension facts are standard prerequisites; the current multi-operation reductions remain fully reproduced.
\par\noindent The source's finite-set working assumption and its infinite-set grading footnote are retained. In particular the source uses ``right-leaning'' within the proof of Proposition 3.8 after defining ``left-leaning'', and its Proposition 3.4 proof uses the symbols $(f_0,x_0)$ and $(x,m_1)$. These are retained as printed, not silently repaired or independently certified.

\setcounter{section}{1}
\setcounter{footnote}{2}

\section{General multi-indices}\label{sec::2}

In applications of the algebraic formalism developed in this paper, one deals with stochastic PDEs of the form
\begin{equation}\label{set01}
\left(\partial_t - \mcL \right) u = \sum_{\mfl\in\mfL^-} a^\mfl (\bu) \xi_\mfl,
\end{equation}
where $u$ is a function of $d+1$ variables $t=x_0$, $x_1$, \dots, $x_d$, $\mcL$ is a differential operator in the variables $x_1$, \dots, $x_d$, $\mfL^- $ is a finite set and the $ \xi_{\mfl} $ are space-time noises (random distributions). The bold letter $\bu$ in each $a^\mfl(\bu)$ means that these are functions of $u$ and of its iterated partial derivatives. For $\bn=(n_0,\dots,n_d)\in\bbN^{d+1}$, one may consider the expressions
\[
u^{(\bn)}:=\frac{\partial_{x_0}^{n_0}\dots \partial_{x_d}^{n_d}}{n_0!\dots n_d!}(u),
\]
which are the coefficients of the Taylor expansion of $u$ around the point $x$. To each $a^\mfl (\bu)$ one may apply the derivatives $\partial_{u^{(\bn)}}$ which pairwise commute. In this context, one typically works with polynomials in the expressions
\[
\left(\prod_{\bn\in\bbN^{d+1}} \partial_{u^{(\bn)}}^{w_{\mfl}(\bn)}\right) (a^\mfl (\bu)),
\]
where $w_{\mfl}(\bn)\in\bbN$ are the orders of derivatives applied to $a^\mfl (\bu)$. These expressions will be now encoded in a formal algebraic way via what we shall call the general multi-indices. For a set $I$, we define the set of multi-indices over $I$, denoted by $M(I)$, as the set of all finitely supported maps $w\colon I\to\bbN$, that is the maps for which all but finitely many $i\in I$ are mapped to $0$. We shall now consider $M(\bbN^{d+1})$, multi-indices over $\bbN^{d+1}$, which one can regard as commutative monomials in letters $\bn\in\bbN^{d+1}$.

Let us now introduce formal variables $z_{(\mfl,w)}$, $(\mfl,w) \in \mfL^-\times M(\bbN^{d+1})$, and define the \emph{general multi-indices} $\beta$ as monomials
\begin{equation}\label{general_multi_indices}
z^{\beta} : = \prod_{(\mfl,w) \in \mfL^-\times M(\bbN^{d+1})} z_{(\mfl,w)}^{\beta(\mfl,w)}.
\end{equation}
These multi-indices have been introduced in the recent work~\cite{BL23}. The vector space spanned by all monomials representing general multi-indices is just the polynomial algebra in variables $z_{(\mfl,w)}$, $(\mfl,w) \in \mfL^-\times M(\bbN^{d+1})$. That algebra has, for each $\bn\in\bbN^{d+1}$, the derivation $D^{(\bn)}$ defined by the formula
\begin{equation}\label{derivations}
D^{(\bn)}=\sum_{(\mfl,w) \in \mfL^-\times M(\bbN^{d+1})} (w(\bn) +1) \, z_{(\mfl,w+e_{\bn})}\partial_{z_{(\mfl,w)}}.
\end{equation}
On the original level of polynomials in iterated partial derivatives of $a^\mfl (\bu)$, this derivation corresponds to applying one extra derivative $\partial_{u^{(\bn)}}$; as in the formula for the derivation $D$ in the introduction, the coefficient $w(\bn)+1$ here corresponds to one of the possible choices of a basis.

We shall be only interested in meaningful multi-indices called \emph{populated} satisfying the condition
\begin{equation}\label{populated_bis}
\sum_{(\mfl,w)} (1 - |w|))\beta(\mfl,w) = 1
\end{equation}
where $ |w|$ is the degree of $w$, viewed as a commutative monomial in letters $\bn\in\bbN^{d+1}$. The vector space spanned by all populated indices is denoted $\mcM^{\genn}$. We can define products $\triangleright_\bn$, $\bn\in\bbN^{d+1}$, on $\mcM^{\genn}$ by setting
\begin{equation}\label{def_products}
z^{\beta} \triangleright_\bn z^{\beta'} = z^{\beta} D^{(\bn)} (z^{\beta'}).
\end{equation}
Our goal now is to have a better understanding of this family of products.

\begin{prop}\label{multi_Nov_example}
Let $A$ be a set and $R$ be a commutative associative algebra, and suppose that $\partial_a$, $a\in A$, are pairwise commuting derivations of that algebra. The bilinear operations $\triangleright_a$ on $R$ defined by the formulas
\[
x\triangleright_a y:=x\partial_a(y)
\]
satisfy the identities
\begin{gather*}
(x \triangleright_a y) \triangleright_b z - x \triangleright_a (y \triangleright_b z) = (y \triangleright_a x) \triangleright_b z - y \triangleright_a (x \triangleright_b z),
\\ (x \triangleright_a y) \triangleright_b z - x \triangleright_a (y \triangleright_b z) = (x\triangleright_b y) \triangleright_a z - x \triangleright_b (y \triangleright_a z),
\\ (x \triangleright_a y) \triangleright_b z = (x \triangleright_b z)\triangleright_a y,
\end{gather*}
for all $ a,b\in A$.
\end{prop}

\begin{proof}
Note that
\[
(x \triangleright_a y) \triangleright_b z=x\partial_a(y)\partial_b(z)=x\partial_b(z)\partial_a(y)=(x \triangleright_b z)\triangleright_a y,
\]
so the last identity holds. Moreover,
\[
x \triangleright_a (y \triangleright_b z)=x\partial_a(y\partial_b(z))=x\partial_a(y)\partial_b(z)+xy\partial_a\partial_b(z),
\]
so
\[
(x \triangleright_a y) \triangleright_b z - x \triangleright_a (y \triangleright_b z)= -xy\partial_a\partial_b(z),
\]
which is symmetric in $x,y$ (since $R$ is commutative) and in $a,b$ (since the derivations $\partial_a$ commute between themselves), proving the first two identities.
\end{proof}

If the set $A$ consists of one element $a$, these identities reduce to the identities
\begin{gather*}
(x \triangleright_a y) \triangleright_a z - x \triangleright_a (y \triangleright_a z) = (y \triangleright_a x) \triangleright_a z - y \triangleright_a (x \triangleright_a z),
\\ (x \triangleright_a y) \triangleright_a z = (x \triangleright_a z)\triangleright_a y,
\end{gather*}
which are the identities of the so-called (left) Novikov algebras~\cite{BN85,DL,O92}. It is thus natural to give the following definition.

\begin{defi}\label{def_multi_novikov}
Let $A$ be a set. A \emph{(left) $A$-multi-Novikov algebra} is a vector space equipped with bilinear products $x,y\mapsto x\triangleright_a y$, $a\in A$, satisfying the identities
\begin{gather}
(x \triangleright_a y) \triangleright_b z - x \triangleright_a (y \triangleright_b z) = (y \triangleright_a x) \triangleright_b z - y \triangleright_a (x \triangleright_b z)\label{eq:exchgen}
,
\\ (x \triangleright_a y) \triangleright_b z - x \triangleright_a (y \triangleright_b z) = (x\triangleright_b y) \triangleright_a z - x \triangleright_b (y \triangleright_a z)\label{eq:exchA}
,
\\ (x \triangleright_a y) \triangleright_b z = (x \triangleright_b z)\triangleright_a y\label{eq:rcom}
,
\end{gather}
for all $ a,b\in A$.
\end{defi}

Note that the first two of the multi-Novikov identities imply the identity
\[
(x \triangleright_a y) \triangleright_b z - x \triangleright_a (y \triangleright_b z) = (y \triangleright_b x) \triangleright_a z - y \triangleright_b (x \triangleright_a z),
\]
which is the defining identity of left multi-pre-Lie algebras~\cite[Proposition~4.21]{BCCH}.

\section{Free multi-Novikov algebras}\label{sec::3}
The only source of multi-Novikov algebras that we have seen so far is given by Proposition~\ref{multi_Nov_example}, which suggests to consider the following algebraic structure.

\begin{defi}
Let $A$ be a set. A commutative multidifferential algebra is a~vector space $R$ equipped with a commutative associative product and endomorphisms $\partial_a$, $a\in A$, which are pairwise commuting derivations of the product on~$R$.
\end{defi}

Note that it is more common to use this structure to define a single binary operation on the different vector space $R\otimes\Vect(A)$, see, for instance~\cite[Section~2.1]{Bur};\linebreak in our work, we use it to define many operations on the same vector space $R$.

Using this terminology, we may re-state Proposition~\ref{multi_Nov_example} by saying that every commutative multidifferential algebra has a natural multi-Novikov structure, which we shall call the \emph{canonical multi-Novikov structure}.

The following almost obvious proposition describes the free commutative multidifferential algebra generated by a set $X$, which we shall denote $\CDi_A\langle X\rangle$. For convenience of the reader, we recall that the notation $M(-)$ refers to multi-indices, that is finitely supported maps to $\bbN$.

\begin{prop}
The algebra $\CDi_A\langle X\rangle$ is isomorphic to the polynomial\linebreak algebra on the set of generators $X\times M(A)$ equipped with the derivations $\partial_a$ defined on generators as follows: for $(x,m)\in X\times M(A)$, we have
\[
\partial_a(x,m)=(x,m_a),
\]
where $m_a(a')=m(a')+\delta_{a,a'}$. In other words, the generator $(x,m)$ of the polynomial algebra is identified with the element $\prod_{a\in A}\partial_a^{m(a)}(x)$ of $\CDi_A\langle X\rangle$.
\end{prop}

The main result of this section is the following theorem that describes the free multi-Novikov algebra generated by a set $X$, which we shall denote $\Nov_A\langle X\rangle$. This theorem is a generalisation of~\cite[Theorem~7.8]{DL}, and follows a similar strategy; however, given how intricate the proof is, we deemed important to give it with all the details.


\begin{theo}\label{th:freeNovikov}
The algebra $\Nov_A\langle X\rangle$ is isomorphic to the multi-Novikov subalgebra of the free commutative multidifferential algebra $\CDi_A\langle X\rangle$ generated by $X$; here we consider $\CDi_A\langle X\rangle$ with its canonical multi-Novikov structure.
\end{theo}

The proof of this theorem consists of several parts. First, let us describe the multi-Novikov subalgebra of the free commutative multidifferential algebra $\CDi_A\langle X\rangle$ generated by $X$ in a more concrete way.

\begin{prop}\label{prop:lowerbound}
The multi-Novikov subalgebra of the free commutative multidifferential algebra $\CDi_A\langle X\rangle$ generated by $X$ is spanned by all monomials $\prod_{x,m}(x,m)^{n_{x,m}}$ for which the ``populated condition'' holds:
\[
\sum_{x,m} n_{x,m}\left(-1+\sum_{a\in A}m(a)\right)=-1.
\]
\end{prop}

\begin{proof}
We shall freely use the identification of generators $x\in X$ with $(x,0)$ in $\CDi_A\langle X\rangle$. Note that the populated condition holds for each generator $x\in X$: we have $m(a)=0$ for all $a\in A$ and $n_{x,m}=1$. Moreover, if we have two monomials $u,v$ for which the populated condition holds, then $u\triangleright_a v=u\partial_a(v)$ is a sum of monomials for which the populated condition holds as well, so the vector space spanned by such monomials is a multi-Novikov subalgebra. Let us show that this subalgebra is generated by~$X$. Let us take some monomial $u$ satisfying the populated condition. If for some $a\in A$ and some $x\in X$ we have $u=\partial_a(x)u_1$, we have $u=u_1\triangleright_a x$, and we may argue by induction on degree. Otherwise, $u$ is a product of a monomial in variables $X$ and a monomial in variables $(x,m)$ where $\sum_{a\in A}m(a)\ge 2$. Let us take one of the latter, say $(x_0,m_0)$ with $\sum_{a\in A}m_0(a)=d\ge 2$, and take some $a_0$ for which $m_0(a_0)\ne 0$. Because of the populated condition, the monomial $u$ is also divisible by a monomial of degree at least $d$ in variables from~$X$, say $u=x_{i_1}\dots x_{i_d}(f_0,x_0)u_1$. Note that the populated condition is satisfied for~$u_1$. Let us define $m_1\in M(A)$ by the formula $m_1(a)=m_0(a)-\delta_{a_0,a}$. Note that the populated condition is satisfied for $x_{i_1}\dots x_{i_d}(x,m_1)$, and we have
\begin{multline}
u_1\triangleright_{a_0}(x_{i_1}\dots x_{i_d}(x,m_1))=\partial_{a_0}(x_{i_1}\dots x_{i_d}(x,m_1))u_1=\\ u+\sum_{j=1}^d (x_{i_1}\dots x_{i_{j-1}} x_{i_{j+1}}\dots x_{i_d}(x,m_1)u_1)\triangleright_{a_0}x_{i_j}
\end{multline}
and we once again may argue by induction on degree.
\end{proof}

We shall use Proposition~\ref{prop:lowerbound} to give a lower bound on the ``size'' of the free multi-Novikov algebra. To make sense of it, we define a $\bbZ\times \bbN$-grading on $\CDi_A\langle X\rangle$ by postulating that each variable $(x,m)$ has the grading
\[
\left(-1+\sum_{a\in A}m(a),1\right),
\]
and extending this grading to products of variables additively. If we assume the sets $A$ and $X$ finite, the bi-homogeneous components of $\CDi_A\langle X\rangle$ are finite-dimensional,\footnote{We shall be making this assumption from now on; for the interested reader, we indicate that in the case where (say) $X$ is infinite, one may consider a finer grading by $\bbN^X$ instead of $\bbN$, which allows for all our proofs to proceed without any problem.} and Proposition~\ref{prop:lowerbound} immediately implies the following lower bound.

\begin{corollary}\label{lower_bound}
For each $n\in\bbN$ the dimension of the degree $n$ component of $\Nov_A\langle X\rangle$ is greater than or equal to the dimension of the $(-1,n)$-homogeneous component of $\CDi_A\langle X\rangle$.
\end{corollary}

We shall now obtain the same upper bound for the dimension of the degree $n$ component of $\Nov_A\langle X\rangle$, which we shall use to prove Theorem~\ref{th:freeNovikov}. We shall for the moment perform all calculations in the multi-magmatic algebra $\Mag_A\langle X\rangle$ which has bilinear operations $\triangleright_a$, $a\in A$, without any identities between them. This algebra has an obvious basis of monomials which are binary trees whose internal vertices are labelled by elements of $A$ and whose leaves are labelled by elements of $X$.

The algebra $\Mag_A\langle X\rangle$ has an ideal $I$ generated by all elements
\begin{gather}
(x \triangleright_a y) \triangleright_b z - x \triangleright_a (y \triangleright_b z) - (y \triangleright_a x) \triangleright_b z + y \triangleright_a (x \triangleright_b z),
\\ (x \triangleright_a y) \triangleright_b z - x \triangleright_a (y \triangleright_b z) - (x\triangleright_b y) \triangleright_a z + x \triangleright_b (y \triangleright_a z),
\\ (x \triangleright_a y) \triangleright_b z - (x \triangleright_b z)\triangleright_a y,
\end{gather}
for all $ a,b\in A$ and all $x,y,z\in \Mag_A\langle X\rangle$. The quotient by this ideal is clearly the free multi-Novikov algebra $\Nov_A\langle X\rangle$, and we shall now obtain a set of elements that spans the quotient $\Mag_A\langle X\rangle/I$.

Let $y_1,\dots,y_{n+1}\in \Mag_A\langle X\rangle$ be some monomials. Recall that monomials $y_1\triangleright_{a_1} (y_2\triangleright_{a_2} (\dots (y_n\triangleright_{a_n} y_{n+1}) \dots))$ are often referred to as right-normed products of $y_1,\dots,y_{n+1}$. Each monomial in $\Mag_A\langle X\rangle/I$ can be uniquely written as a right-normed product $y_1\triangleright_{a_1} (y_2\triangleright_{a_2} (\dots (y_n\triangleright_{a_n} z) \dots))$, where $z\in X$; for typographic purposes, we shall denote such product as $r(y_1,a_1,\dots,y_n,a_n;z)$.

\begin{defi}\label{leaning}
A right-normed product $r(y_1,a_1,\dots,y_n,a_n;z)$ with $z\in X$ is said to be \emph{left-leaning} if $y_2,\dots,y_n\in X$, and $y_{1}$ is either a generator or a monomial whose right-normed product decomposition is itself left-leaning.
\end{defi}

Let us furthermore specify a subclass of left-leaning products.

\begin{defi}\label{ordered}
Let us fix an order of the set $A$ and an order of the set $X$. We shall say that a left-leaning product $r(y_1,a_1,\dots,y_n,a_n;z)$ is \emph{ordered} if $n=0$ (so that the whole product is just the generator $z$) or $n>0$, $a_1\le a_2\le\dots\le a_n$, and\pagebreak{} $y_1\in X$ or $y_1=r(y_1',a_1',\dots,y_m',a_m';z')$ is an ordered left-leaning product satisfying one of the two conditions:
\begin{itemize}
\item $n<m$,
\item $n=m$ and $z<z'$.
\end{itemize}
\end{defi}

\begin{prop}\label{prop:ordright}
The quotient $\Mag_A\langle X\rangle/I$ is spanned by cosets of ordered left-leaning products.
\end{prop}

\begin{proof}
Let us prove that in $\Mag_A\langle X\rangle/I$ every coset of a monomial $u$ of degree $d$ can be written as a linear combination of cosets of monomials of degree $d$ which are ordered left-leaning products. We shall argue by induction on $d$, and for any fixed $d$ by induction on $n$, the length of the right-normed product decomposition $r(y_1,a_1,\dots,y_n,a_n;z)$ of our monomial~$u$.

Clearly, Identity~\eqref{eq:exchA} implies that we can put $a_1,\dots,a_n$ in the right order at the cost of elements with same $d$ and smaller $n$ that appear (to which the induction hypothesis applies). Similarly, Identity~\eqref{eq:exchgen} implies that we can put $y_1,\dots,y_{n}$ in the decreasing order according to their degrees at the cost of elements with same $d$ and smaller $n$ that appear (to which the induction hypothesis applies). If the resulting element
\[
u'=r\left(z_1,b_1,\dots,z_n,b_n;z\right)
\]
satisfies $z_1\in X$, then the fact that the degrees of $z_1,\dots,z_{n}$ are in the decreasing order implies that we have $z_i\in X$ for all $i$, and so $u'$ is an ordered right-leaning product. If $z_1\notin X$, then the induction hypothesis applies to it, and its coset contains a linear combination of ordered right-leaning products. Without loss of generality, we may take
\[
z_1=r\left(z_1',b_1',\dots,z_m',b_m';z'\right)
\]
with $z_2',\dots,z_m'\in X$. Using Identity~\eqref{eq:rcom} for $a=b_1$ and $b=b_1'$, we see that in the quotient $\Mag_A\langle X\rangle/I$ we have
\[
u'=r\left(r\left(z_1',b_1,z_2,\dots,z_n,b_n;z\right),b_1',z_2',\dots,z_m',b_m';z'\right)
\]
Now we can use the induction hypothesis to write the monomial
\[
r(z_1',b_1,z_2,\dots,z_n,b_n;z)
\]
as a linear combination of ordered right-leaning products. To ensure the same for the monomial $u'$, we may have to use Identity~\eqref{eq:rcom} once again, completing the proof.
\end{proof}

We are now ready to describe what will be a basis of the free multi-Novikov algebra. For that, we shall first recall and use some standard combinatorial argument (Knuth's rotation correspondence~\cite[Chapter~2.3]{Kn}) relating planar rooted binary trees and planar rooted trees. Let $T$ be a planar rooted binary tree with $n$ leaves. We define a planar rooted tree $\rho(T)$ with $n$ vertices as follows: for each internal vertex $i$ of $T$, consider the path going from $i$ to the rightmost leaf of the tree growing from $i$, and collapse that path to a single vertex. An example of this is displayed in Figure~\ref{fig1}, where we use the same label for all the vertices of the planar binary tree collapsed to the same one by the map $\rho$.
\begin{figure}[h]
\centering
\begin{tikzcd}[sep=small]
&&& 6 && 2 \\
5 && 4 && 2 \\
& 4 &&&&&&&& 5 \\
1 && 2 &&&&&& 1 & 4 & 6 \\
& 2 && 3 &&&&&& 2 \\
&& 3 &&&&&&&& 3
\arrow[no head, from=1-4, to=2-5]
\arrow[no head, from=2-1, to=3-2]
\arrow[no head, from=2-5, to=1-6]
\arrow[no head, from=3-2, to=2-3]
\arrow[no head, from=3-2, to=4-3]
\arrow[no head, from=3-10, to=4-10]
\arrow[no head, from=4-3, to=2-5]
\arrow[no head, from=4-10, to=5-10]
\arrow[no head, from=4-11, to=5-10]
\arrow[no head, from=5-2, to=4-1]
\arrow[no head, from=5-2, to=4-3]
\arrow[no head, from=5-10, to=4-9]
\arrow[no head, from=6-3, to=5-2]
\arrow[no head, from=6-3, to=5-4]
\arrow[no head, from=6-11, to=5-10]
\end{tikzcd}
\caption{Knuth's rotation correspondence}\label{fig1}
\end{figure}

\goodbreak
If $T$ is the underlying planar binary tree of a monomial $u\in \Mag_A\langle X\rangle$, we have a~labelling of all internal vertices of $T$ by elements of the set $A$ and a labelling of all leaves of $T$ by elements of the set $X$. These labellings induce the labelling of vertices and edges of the planar tree $\rho(T)$: each vertex of $\rho(T)$ acquires the label of the leaf of $T$ which we collapsed onto that vertex, and each edge of $\rho(T)$ acquires the label of the vertex of $T$ for which that edge is the left incoming edge. An example of this is displayed in Figure~\ref{fig2}.

\begin{figure}[h]
\centering
\begin{tikzcd}[sep=small]
&&& x && x \\
y && x && a \\
& b &&&&&&&& y \\
x && b &&&&&& x & x & x \\
& a && y &&&&&& x \\
&& b &&&&&&&& y
\arrow[no head, from=1-4, to=2-5]
\arrow[no head, from=2-1, to=3-2]
\arrow[no head, from=2-5, to=1-6]
\arrow[no head, from=3-2, to=2-3]
\arrow[no head, from=3-2, to=4-3]
\arrow["b"', no head, from=3-10, to=4-10]
\arrow[no head, from=4-3, to=2-5]
\arrow["b"', no head, from=4-10, to=5-10]
\arrow["a"', no head, from=4-11, to=5-10]
\arrow[no head, from=5-2, to=4-1]
\arrow[no head, from=5-2, to=4-3]
\arrow["a", no head, from=5-10, to=4-9]
\arrow[no head, from=6-3, to=5-2]
\arrow[no head, from=6-3, to=5-4]
\arrow["b", no head, from=6-11, to=5-10]
\end{tikzcd}
\caption{Knuth's rotation correspondence with labels}\label{fig2}
\end{figure}

We shall refer to an occurrence of a generator $x\in X$ in $u$ as a
\emph{root} if it becomes a~label of an internal vertex of $\rho(T)$, and a
\emph{leaf} if it becomes a label of a leaf of $\rho(T)$.

\begin{defi}\label{normal}
We shall call an ordered left-leaning product \emph{normal} if the sequence of its leaves, read from the left to the right, is in the increasing order with respect to the order of $X$.
\end{defi}

We are now ready to prove the main technical result of this section.

\begin{prop}\label{normal_coset}
The quotient $\Mag_A\langle X\rangle/I$ is spanned by cosets of normal ordered left-leaning products.
\end{prop}


\begin{proof}
Thanks to the result of Proposition~\ref{prop:ordright}, it is enough to prove that every ordered left-leaning product can be written as a linear combination of cosets of normal ordered left-leaning products. Let us take an ordered left-leaning product
\[
u=r(y_1,a_1,\dots,y_n,a_n;z)
\]
with $a_1\le a_2\le\dots\le a_n$, $y_2,\dots,y_n\in X$, and either $y_{1}\in X$ or
\[
y_{1}=
r\left(y_1',a_1',\dots,y_m',a_m';z'\right)
\]
satisfies $n<m$ or $n=m$ and $z<z'$. We shall prove that in $\Mag_A\langle X\rangle/I$ it can be written as a linear combination of cosets of normal ordered left-leaning products by induction on $d$, and for a fixed $d$ by induction on $\ell$, the number of leaves, and for fixed $d$ and $\ell$ by the sum $n+m$.

Suppose first that $y_{1}\in X$, so that all the elements $y_1,\dots,y_{n}$ are leaves. Using Identity~\eqref{eq:exchgen}, we may put them in order at the cost of elements with fewer leaves that appear. We may apply the procedure of Proposition~\ref{prop:ordright} (which does not increase the number of leaves) to these elements, and then the induction hypothesis finishes the argument.

Suppose now that $y_1\notin X$, so that we have $y_{1}=r(y_1',a_1',\dots,y_m',a_m';z')$ with $n<m$ or $n=m$ and $z<z'$. By the induction hypothesis we may assume that $y_{1}$ is a normal ordered left-leaning product. Also, modulo terms with fewer leaves, we may assume $y_2\le y_3\le\dots\le y_n$. Suppose $y_m'> y_2$. If we show that these two elements can be exchanged modulo $I$, the proof will be finished after doing several exchanges like that.

Using our relations, we can change the order from $r(y_1',a_1',\dots,y_m',a_m';z')$ to $r(y_m',a_1',y_1'\dots,y_{m-1}',a_m';z')$ at the cost of some elements with fewer leaves, and one more element $r((y_1'\triangleright a_1' y_m'),a_2,\dots,y_{m-1}',a_m';z')$ with the same number of leaves. However, this element has fewer factors in the second level right-normed factorisation, and the induction on $n+m$ takes care of that. Thus, we may focus on the element
\[
r\left(r(y_m',a_1',y_1'\dots,y_{m-1}',a_m';z'),a_1,y_2,\dots,y_n,a_n;z\right),
\]
which, thanks to Identity~\eqref{eq:rcom}, is equal to
\[
r\left(r(y_m',a_1,y_2,\dots,y_n,a_n;z),a_1',y_1'\dots,y_{m-1}',a_m';z'
\right),
\]
Now, as before, we use Identity~\eqref{eq:exchgen} to exchange in $r(y_m',a_1,y_2,\dots,y_n,a_n;z)$ the elements $y_m'$ and $y_2$, and use Identity~\eqref{eq:rcom} again to obtain
\[
r\left(r(y_2,a_1',y_1'\dots,y_{m-1}',a_m';z'),a_1,y_m',\dots,y_n,a_n;z
\right),
\]
and an argument similar to before allows us to exchange $y_2$ and $y_{1}'$ in the inner monomial $r(y_2,a_1',y_1',\dots,y_{m-1}',a_m';z')$, leading to a linear combination of normal ordered left-leaning products by the induction hypothesis.
\end{proof}

\begin{proof}[{Proof of Theorem~\ref{th:freeNovikov}}]
Let us show that normal ordered left-leaning products of degree $d$ are in bijection with the standard basis elements of the $(-1,d)$-homogeneous component of $\CDi_A\langle X\rangle$. The way to construct a basis element out of an ordered left-leaning product is quite straightforward: to each root vertex labelled by $x$ whose operations in the right-normed factorisation are $\triangleright_{a_1},\dots, \triangleright_{a_n}$ with $a_1\le a_2\le a_n$, we associate the variable $\partial_{a_1}\dots\partial_{a_n}(x)$, and to each leaf labelled by $x$ we associate simply the variable $x$. The reverse reconstruction is straightforward (the ordered condition of the left-leaning product ensures uniqueness), and leads to a valid monomial in $\Mag_A\langle X\rangle$ thanks to the populated condition.

Since dimensions of bi-homogeneous components of $\CDi_A\langle X\rangle$ are finite-dimensional, we see that the surjective map
\[
\Nov_A\langle X\rangle\twoheadto\CDi_A\langle X\rangle_{-1}
\]
must be an isomorphism, as required.
\end{proof}

% Faithful selected bibliography transcription from the official published PDF, printed pp.812--813.
% Native citation keys mapped to the printed labels of that same version; page-backlinks omitted.
\begin{thebibliography}{BCCH21}
\bibitem[BCCH21]{BCCH} Yvain Bruned, Ajay Chandra, Ilya Chevyrev, and Martin Hairer, \emph{Renormalising SPDEs in regularity structures}, J. Eur. Math. Soc. 23 (2021), no. 3, 869--947.
\bibitem[BL24]{BL23} Yvain Bruned and Pablo Linares, \emph{A top-down approach to algebraic renormalization in regularity structures based on multi-indices}, Arch. Ration. Mech. Anal. 248 (2024), no. 6, article no. 111 (81 pages).
\bibitem[BN85]{BN85} Alexander A. Balinski\u\i{} and Serge\u\i{} P. Novikov, \emph{Poisson brackets of hydrodynamic type, Frobenius algebras and Lie algebras}, Dokl. Akad. Nauk SSSR 283 (1985), no. 5, 1036--1039.
\bibitem[Bur06]{Bur} Dietrich Burde, \emph{Left-symmetric algebras, or pre-Lie algebras in geometry and physics}, Cent. Eur. J. Math. 4 (2006), no. 3, 323--357.
\bibitem[DL02]{DL} Askar Dzhumadil'daev and Clas L\"ofwall, \emph{Trees, free right-symmetric algebras, free Novikov algebras and identities}, Homology Homotopy Appl. 4 (2002), no. 2, 165--190.
\bibitem[Knu97]{Kn} Donald E. Knuth, \emph{The art of computer programming. Vol. 1. Fundamental Algorithms}, 3rd ed., Addison-Wesley Publishing Group, 1997.
\bibitem[Osb92]{O92} J. Marshall Osborn, \emph{Novikov algebras}, Nova J. Algebra Geom. 1 (1992), no. 1, 1--13.
\end{thebibliography}
\end{document}
